(1) 恒温共轭张力为 \(f=(\partial A/\partial L)_T=Nk_BT(L-L_0)/L_0^2\)，正拉伸时 \(L>L_0\) 有 \(f>0\)。熵 \(S=-(\partial A/\partial T)_L=S_0(T)-Nk_B(L-L_0)^2/(2L_0^2)\)，表明限制高分子链构型使熵减少。
(2) 等温可逆拉伸的机械功 \(W=\int_{L_i}^{L_f}f\,dL=Nk_BT[(L_f-L_0)^2-(L_i-L_0)^2]/(2L_0^2)=\Delta A\)。温度不变时 \(\Delta S=-Nk_B[(L_f-L_0)^2-(L_i-L_0)^2]/(2L_0^2)\)，故系统吸热 \(Q=T\Delta S=-W\)；若远离无应力长度的平方伸长增加，元件放热，若恢复向无应力长度则吸热。
(3) 绝热可逆路径满足 \(\Delta S=0\)，于是 \(C\ln(T_f/T_i)=Nk_B[(L_f-L_0)^2-(L_i-L_0)^2]/(2L_0^2)\)。因此 \(T_f=T_i\exp\{Nk_B[(L_f-L_0)^2-(L_i-L_0)^2]/(2CL_0^2)\}\)；增加链的拉伸量平方使绝热温度上升，该结论依赖理想熵弹簧及常热容假设。
(4) \((\partial f/\partial T)_L=Nk_B(L-L_0)/L_0^2\)，而 \(-(\partial S/\partial L)_T=Nk_B(L-L_0)/L_0^2\)，两者一致。若 \(L_f=L_i\) 则 \(W=Q=0,T_f=T_i\)；若 \(N\to0\)，链构型贡献消失，所给弹性模型无张力、无热弹耦合。
